%\section*{Lernziele: Excel}
%\begin{todolist}
%\item Ich verstehe das Koordinatensystem von Excel-Zellen (z.B. $B4$ für eine Zelle in Kolonne 2, Zeile 4)
%\item Ich kann Zellen als Zahlen, Text, Währung, Datum und in weiteren Formaten darstellen.
%\item Ich kann einfach Formeln erstellen, beispielsweise um den Mittelwert, die Summe oder das Maximum einer Datenreihe zu bestimmen.
%\item Ich verstehe den Unterschied zwischen fixierten und nicht-fixierten Zellreferenzen (z.B. $B2$ vs. $\$B\$2$) und kann fixierte Zellreferenzen verwenden, um Formeln zu kopieren.
%\item Ich kann Datenreihen als Punkt-, Kreis- oder Liniendiagramm darstellen und kann, je nach Kontext, passende, weitere Diagrammelemente hinzufügen, wie beispielsweise eine Legende, X-Achsen-Titel, Y-Achsen-Titel, Grafik-Titel oder Führungslinien. Ich kann zudem den Bereich er X- und Y-Achse manuell anpassen.
%\end{todolist}


\chapter{Lernziele}
\section{Korrelation und Kausalität}
\begin{todolist}
    \item Ich kann den Unterschied zwischen Kausalität und Korrelation anhand von konkreten Beispielen erläutern.
    \item Ich kann das gemeinsame Verhältnis zwischen zwei Variablen A und B inhaltlich korrekt bestimmen, indem ich unterscheide folgende Fälle unterscheide:
    \begin{todolist}
        \item Direkte Kausalität ($A\rightarrow B$)
        \item Indirekte Kausalität ($A\rightarrow C\rightarrow B$)
        \item Gemeinsamer Grund ($C\rightarrow A$ und $C\rightarrow B$)
        \item Zyklische Beeinflussung ($A\rightarrow B \rightarrow A$)
        \item Koinzidenz (= zufällige Korrelation)
    \end{todolist}
    \item Ich weiss, was positive bzw. negative Korrelation ist und kann Beispiele für beide Fälle nennen.
    \item Ich kann anhand von Punktdiagrammen erkennen, ob zwei Grössen positiv,
    negativ oder gar nicht korrelieren (sofern dies einigermassen deutlich zu
    erkennen ist).
    \item Ich verstehe den Ausdruck $Z=x_1'\cdot y_1'+x_2'\cdot y_2'+...+x_n'\cdot y_n'$.
    \item Ich kann den Unterschied zwischen $Z$ und dem Korrelationskoeffizienten $r$ erklären.
    \item Ich kann den Korrelationskoeffizienten in einfachen Beispielen selbst berechnen
    (mit TR).
    \item Ich weiss, dass der Korrelationskoeffizient einheitenlos ist.
    \item Ich weiss, dass der Korrelationskoeffizient immer zwischen -1 und +1 liegt.
\end{todolist}


\section{Einführung}
\begin{todolist}
    \item Ich verstehe, was ein Algorithmus ist und kann ihn von einem Programm unterscheiden.
    \item Ich kenne den Unterschied zwischen traditioneller Programmierung und Maschinellem Lernen.
    \item Ich kann \textbf{Klassifikation} und \textbf{Regression} anhand von Beispielen unterscheiden.
    \item Ich kenne die vier Funktionsarten von Algorithmen: Priorisierung, Klassifikation, Assoziation und Filterung.
    \item Ich verstehe die Unterschiede zwischen \textbf{regelbasierten} und \textbf{selbstlernenden} Algorithmen.
\end{todolist}

\section{Datenvorbereitung und Evaluation}
\begin{todolist}
    \item Ich verstehe, warum ein \textbf{Train-Test-Split} notwendig ist und kann den Unterschied zwischen Trainings- und Testdaten erklären.
    \item Ich kann \textbf{Features} (Merkmale) und \textbf{Label} (Zielvariable) in einem Datensatz identifizieren.
    \item Ich kann die \textbf{Min-Max-Normalisierung} durchführen und verstehe, warum sie bei distanzbasierten Algorithmen wichtig ist.
    \item Ich kenne die \textbf{Konfusionsmatrix} und kann daraus \textbf{Accuracy}, \textbf{Precision} und \textbf{Recall} berechnen.
    \item Ich kann anhand dieser Metriken beurteilen, ob ein Klassifikationsmodell gut funktioniert.
\end{todolist}

\section{Priorisierung}
\begin{todolist}
    \item Ich verstehe die Grundidee von \textbf{Priorisierungsalgorithmen} und kann sie von Klassifikation unterscheiden.
    \item Ich kann \textbf{gewichtete Ranglisten} erstellen und die Gewichte begründen.
    \item Ich verstehe die Grundidee von \textbf{PageRank} und kann eine einfache Berechnung für eine Iteration durchführen.
    \item Ich verstehe die Risiken von Priorisierungsalgorithmen, insbesondere Manipulation und Bias-Verstärkung, und kann Beispiele dafür nennen.
\end{todolist}

\section{Entscheidungsbaum}
\begin{todolist}
    \item Ich verstehe die Grundidee eines \textbf{Entscheidungsbaums} und kann einen Baum lesen.
    \item Ich kann den \textbf{Gini-Koeffizienten} berechnen und interpretieren.
    \item Ich verstehe das Problem des \textbf{Overfitting} und kann es anhand von Trainings- und Testgenauigkeit erkennen.
    \item Ich kann Hyperparameter wie \texttt{max\_depth} nutzen, um Overfitting zu vermeiden.
    \item Ich erkenne, dass die hohe Interpretierbarkeit von Entscheidungsbäumen nicht vor Bias schützt, wenn Trainingsdaten verzerrt sind.
\end{todolist}

\section{Random Forest}
\begin{todolist}
    \item Ich verstehe die Grundidee von \textbf{Random Forests} und die Rolle von \textbf{Bagging} und \textbf{Feature-Subsampling}.
    \item Ich kann \textbf{Feature Importance} interpretieren und verstehe ihre Bedeutung.
    \item Ich kenne die Hyperparameter \texttt{n\_estimators}, \texttt{max\_depth} und \texttt{min\_samples\_leaf} und kann deren Effekt auf die Modellleistung erklären.
    \item Ich verstehe das \textbf{Black-Box-Problem} von Ensembles und erkenne die Risiken der Bias-Verstärkung bei systematischen Verzerrungen in den Trainingsdaten.
    \item Ich kann begründen, warum Interpretierbarkeit in Hochrisiko-Entscheidungen wichtiger sein kann als maximale Accuracy.
\end{todolist}

\section{k-Nearest Neighbors}
\begin{todolist}
    \item Ich verstehe die Grundidee von \textbf{KNN} und wie Klassifikation durch Ähnlichkeit funktioniert.
    \item Ich kann die \textbf{euklidische Distanz} berechnen und verstehe ihre Bedeutung für KNN.
    \item Ich kann den Einfluss von $k$ auf Overfitting und die Glätte der Entscheidungsgrenze erklären.
    \item Ich verstehe, dass Normalisierung bei KNN obligatorisch ist und Wertsetzungen enthält.
    \item Ich erkenne, wie Normalisierungsbias und fehlerhafte Features zu Diskriminierung führen können.
\end{todolist}

\section{Markov-Ketten}
\begin{todolist}
    \item Ich verstehe die \textbf{Markov-Eigenschaft} und kann überprüfen, ob sie in einer gegebenen Situation erfüllt ist.
    \item Ich kann Markov-Ketten als \textbf{gerichtete Graphen} und als \textbf{Übergangsmatrizen} darstellen.
    \item Ich kann Übergangswahrscheinlichkeiten berechnen und Zustände vorhersagen.
    \item Ich kann die Darstellungsformen ineinander überführen.
\end{todolist}


\section{Support Vector Machine (Optional)}
\begin{todolist}
    \item Ich verstehe die Grundidee von \textbf{SVMs}: das Finden einer Trennebene mit maximalem \textbf{Margin}.
    \item Ich kann erklären, warum der Margin wichtig ist und wie der Parameter $C$ den Trade-off zwischen Trainingsgenauigkeit und Generalisierung steuert.
    \item Ich verstehe das Konzept des \textbf{Kernel-Tricks} und kann erklären, wie er nichtlinear trennbare Daten handhabt.
    \item Ich erkenne, dass Kernel-Wahl und andere Hyperparameter Wertsetzungen sind und direkte Auswirkungen auf Fairness haben.
\end{todolist}

\section{Apriori}
\begin{todolist}
    \item Ich verstehe, dass Apriori \textbf{häufige Muster} in Transaktionsdaten findet, keine Vorhersagen trifft.
    \item Ich kann \textbf{Support}, \textbf{Confidence} und \textbf{Lift} berechnen und interpretieren.
    \item Ich verstehe, dass Lift $> 1$ für positive Assoziationen spricht.
    \item Ich erkenne die Gefahr von \textbf{Scheinkorrelationen} und kann begründen, warum Schwellenwerte notwendig sind.
\end{todolist}

\section{Naive Bayes}
\begin{todolist}
    \item Ich verstehe das \textbf{Bayes-Theorem} und kann es anwenden, um Wahrscheinlichkeiten zu berechnen.
    \item Ich kann erklären, was die \enquote{naive} Annahme der bedingten Unabhängigkeit ist.
    \item Ich verstehe, dass diese Annahme oft verletzt ist, aber Naive Bayes trotzdem praktisch funktioniert.
    \item Ich kann \textbf{Laplace-Glättung} anwenden, um Wahrscheinlichkeitsprobleme bei unbekannten Wörtern zu vermeiden.
    \item Ich kann Naive Bayes für Filterung und Klassifikation einsetzen und verstehe die Rolle von Precision und Recall in solchen Kontexten.
\end{todolist}

\section{Regression}
\begin{todolist}
	\item Ich kann eine Gerade zwischen zwei Punkten \textbf{manuell berechnen}.
	\item Ich kann erklären, was man unter der \textbf{abhängigen Variable} versteht: Die abhängige Variable ist die Zielvariable, deren Wert durch die unabhängigen Variablen vorhergesagt werden soll.
	\item Ich kann erklären, was man unter \textbf{unabhängigen Variablen} versteht: Unabhängige Variablen sind die Einflussgrössen, mit deren Hilfe die abhängige Variable modelliert wird.
	\item Ich kenne den Unterschied zwischen \textbf{einfacher} und \textbf{multipler} linearer Regression: Bei der einfachen linearen Regression gibt es nur eine unabhängige Variable, bei der multiplen linearen Regression mehrere.
	\item Ich weiss, weshalb man die \textbf{Linearität der unabhängigen Variablen} überprüfen muss: Die Beziehung zwischen den unabhängigen Variablen und der abhängigen Variable sollte linear sein, da sonst die Ergebnisse der Regression nicht zuverlässig sind. Die Überprüfung erfolgt z.B. durch Streudiagramme oder Korrelationsanalysen.
	\item Ich kann die Qualität einer Regression mit der Summe der \textbf{Residuen} und der \textbf{Summe der quadrierten Residuen} beurteilen: Die Residuen sind die Abweichungen zwischen den tatsächlichen Werten und den vorhergesagten Werten. Die Summe der quadrierten Residuen (RSS) gibt an, wie gut die Regressionsgerade zu den Daten passt; je kleiner die RSS, desto besser die Anpassung.
	\item Ich weiss, wozu man die Programmbibliothek \texttt{pandas} verwendet: \texttt{pandas} dient zur effizienten Datenverwaltung, -analyse und -vorverarbeitung in Python.
	\item Ich weiss, wozu man die Programmbibliothek \texttt{statsmodels} verwendet: \texttt{statsmodels} wird zur Durchführung statistischer Analysen und zur Erstellung von Regressionsmodellen in Python genutzt.
	\item Ich kann erklären, was ein \textbf{statistisches Modell} ist: Ein statistisches Modell ist eine mathematische Beschreibung, die die Beziehung zwischen Variablen in einem Datensatz abbildet und Vorhersagen ermöglicht.
\end{todolist}

\section{Modellvergleich}
\begin{todolist}
    \item Ich kann verschiedene Klassifikationsmodelle (Entscheidungsbaum, Random Forest, SVM, KNN, Naive Bayes) vergleichen und ihre Stärken/Schwächen benennen.
    \item Ich kann lineare Regression mit anderen Regressionsmodellen vergleichen.
    \item Ich kann für ein gegebenes Problem das passende Modell auswählen und begründen.
    \item Ich verstehe den Trade-off zwischen Interpretierbarkeit und Genauigkeit und kann begründen, wann welcher wichtiger ist.
    \item Ich kann \textbf{Datenbias}, \textbf{Modellbias} und deren Auswirkungen erkennen.
    \item Ich kann konkrete Massnahmen zur Bias-Reduktion vorschlagen.
\end{todolist}

\section{Neuronale Netze}
\begin{todolist}
    \item Ich verstehe das Konzept eines \textbf{Perceptron} und kann erklären, warum ein einzelnes Neuron nicht XOR lernen kann.
    \item Ich kann die Architektur eines \textbf{Multilayer Perceptron (MLP)} beschreiben (Eingabe-, versteckte und Ausgabeschichten).
    \item Ich kenne häufige \textbf{Aktivierungsfunktionen} (Sigmoid, ReLU, Softmax) und kann ihre Verwendung erklären.
    \item Ich verstehe das Konzept von \textbf{Backpropagation} und Gradientenabstieg, ohne die mathematischen Details formal beweisen zu müssen.
    \item Ich erkenne, dass neuronale Netze sehr leistungsfähig sind, aber viele Daten und hohe Rechenleistung benötigen.
    \item Ich verstehe das \textbf{Black-Box-Problem} von tiefen Netzen und erkenne, wann Transparenz kritisch ist.
\end{todolist}

\section{Transferaufgabe und Projektarbeit}
\begin{todolist}
    \item Ich kann ein eigenes Datenprojekt von der Problemdefinition über Datenbeschaffung, Vorbereitung, Modellierung und Visualisierung durchführen.
    \item Ich kann mindestens zwei Modelle trainieren, vergleichen und mit geeigneten Metriken bewerten.
    \item Ich kann Bias-Risiken in meinem Projekt erkennen und kritisch reflektieren.
    \item Ich kann meine Ergebnisse in einer 3-Minuten-Präsentation zusammenfassen und begründen.
\end{todolist}

\chapter{Datenanalyse und -visualisierung}
\begin{todolist}
  \item Ich kann Daten aus gängigen Datenformaten (insbesondere \texttt{.csv}) in Python einlesen und analysieren.
  \item Ich kann Datentypen ausgeben und verstehe die Unterschiede zwischen den Datentypen \texttt{int64}, \texttt{float64}, \texttt{object} und \texttt{bool}.
  \item Ich kann einzelne Spalten mittels Befehlen der Library \texttt{pandas} auswählen und anzeigen.
  \item Ich kann grundlegende Aggregationsfunktionen der Library \texttt{pandas} auf Tabellen und Spalten anwenden.
  \item Ich kann mit \texttt{pandas} Daten filtern und spezifische Informationen extrahieren.
  \item Ich kann mit \texttt{pandas} Ergebnisse auf unterschiedliche Weise visualisieren (Kuchen-, Balken-, Streudiagramm, Histogramme) und interpretieren.
\end{todolist}